%%%%%%%%%%%%%%%%%%%%%%%%%%%%%%%%%%%%%%%%%%%%%%%%%%%%%%%%%%%%%%%%%%%%%%%%%%%%%%%%%%
\begin{frame}[fragile]\frametitle{}
\begin{center}
{\Large AI Fluency Framework}

{\tiny (Ref: Anthropic Academy Course)}
\end{center}
\end{frame}

%%%%%%%%%%%%%%%%%%%%%%%%%%%%%%%%%%%%%%%%%%%%%%%%%%%%%%%%%%%%%%%%%%%%%%%%%%%%%%%%%%
\begin{frame}[fragile]\frametitle{What is AI Fluency?}
\begin{columns}
    \begin{column}[T]{0.6\linewidth}
      \begin{itemize}
        \item \textbf{AI Fluency} is the ability to work with AI systems in ways that are:
          \begin{itemize}
            \item Effective: achieving your intended goals
            \item Efficient: making good use of time and resources
            \item Ethical: respecting values and responsibilities
            \item Safe: avoiding harm to people and systems
          \end{itemize}
        \item It combines:
          \begin{itemize}
            \item Practical skills and domain knowledge
            \item Critical insight and contextual judgment
            \item Values that guide responsible AI use
            \item Adaptability to evolving AI technologies
          \end{itemize}
      \end{itemize}
    \end{column}
    \begin{column}[T]{0.4\linewidth}
      \begin{block}{The 4Ds Framework}
        The four core competencies of AI Fluency:
        \begin{enumerate}
          \item \textbf{Delegation}
          \item \textbf{Description}
          \item \textbf{Discernment}
          \item \textbf{Diligence}
        \end{enumerate}
      \end{block}
    \end{column}
  \end{columns}
\end{frame}

%%%%%%%%%%%%%%%%%%%%%%%%%%%%%%%%%%%%%%%%%%%%%%%%%%%%%%%%%%%%%%%%%%%%%%%%%%%%%%%%%%
\begin{frame}[fragile]\frametitle{The 4Ds: Delegation}
\begin{columns}
    \begin{column}[T]{0.6\linewidth}
      \begin{itemize}
        \item \textbf{Delegation}: Deciding what work should be done by humans, by AI, or by both together
        \item Key sub-skills:
          \begin{itemize}
            \item \textbf{Problem Awareness:} Clearly understanding your goals and the nature of the work \emph{before} involving AI
            \item \textbf{Platform Awareness:} Understanding the capabilities and limitations of different AI systems
            \item \textbf{Task Delegation:} Thoughtfully distributing work between humans and AI to leverage the strengths of each
          \end{itemize}
        \item Good delegation means neither under-using nor over-relying on AI
      \end{itemize}
    \end{column}
    \begin{column}[T]{0.4\linewidth}
      \begin{alertblock}{Key Question}
        \emph{``Which parts of this task are best handled by AI, and which require human judgment?''}
      \end{alertblock}
      \vspace{0.5em}
      \begin{exampleblock}{Example}
        Delegating first-draft writing to AI while retaining human oversight for factual accuracy and tone
      \end{exampleblock}
    \end{column}
  \end{columns}
\end{frame}

%%%%%%%%%%%%%%%%%%%%%%%%%%%%%%%%%%%%%%%%%%%%%%%%%%%%%%%%%%%%%%%%%%%%%%%%%%%%%%%%%%
\begin{frame}[fragile]\frametitle{The 4Ds: Description}
\begin{columns}
    \begin{column}[T]{0.6\linewidth}
      \begin{itemize}
        \item \textbf{Description}: Effectively communicating with AI systems
        \item Three dimensions:
          \begin{itemize}
            \item \textbf{Product Description:} Define the desired \emph{output}: format, audience, style, length
            \item \textbf{Process Description:} Define \emph{how} the AI should approach the task: step-by-step instructions, constraints, order of operations
            \item \textbf{Performance Description:} Define the AI's \emph{behaviour} during the interaction: concise vs. detailed, challenging vs. supportive
          \end{itemize}
        \item Clear description reduces guesswork and improves output quality
      \end{itemize}
    \end{column}
    \begin{column}[T]{0.4\linewidth}
      \begin{alertblock}{Key Question}
        \emph{``Have I given the AI enough context to know what I actually need?''}
      \end{alertblock}
      \vspace{0.5em}
      \begin{exampleblock}{Example}
        ``Summarise this report in three bullet points for a non-technical executive audience, using plain language''
      \end{exampleblock}
    \end{column}
  \end{columns}
\end{frame}

%%%%%%%%%%%%%%%%%%%%%%%%%%%%%%%%%%%%%%%%%%%%%%%%%%%%%%%%%%%%%%%%%%%%%%%%%%%%%%%%%%
\begin{frame}[fragile]\frametitle{The 4Ds: Discernment}
\begin{columns}
    \begin{column}[T]{0.6\linewidth}
      \begin{itemize}
        \item \textbf{Discernment}: Thoughtfully and critically evaluating AI outputs, processes, and behaviour
        \item Three dimensions:
          \begin{itemize}
            \item \textbf{Product Discernment:} Evaluate the \emph{quality} of what AI produces: accuracy, appropriateness, coherence, relevance
            \item \textbf{Process Discernment:} Evaluate \emph{how} the AI arrived at its output: check for logical errors, lapses in attention, or flawed reasoning steps
            \item \textbf{Performance Discernment:} Evaluate \emph{how} the AI behaves during interaction: is its communication style effective for your needs?
          \end{itemize}
      \end{itemize}
    \end{column}
    \begin{column}[T]{0.4\linewidth}
      \begin{alertblock}{Key Question}
        \emph{``Is this output accurate, appropriate, and free from subtle errors?''}
      \end{alertblock}
      \vspace{0.5em}
      \begin{exampleblock}{Watch out for}
        \begin{itemize}
          \item Hallucinations
          \item Plausible-sounding but incorrect facts
          \item Bias in generated content
        \end{itemize}
      \end{exampleblock}
    \end{column}
  \end{columns}
\end{frame}

%%%%%%%%%%%%%%%%%%%%%%%%%%%%%%%%%%%%%%%%%%%%%%%%%%%%%%%%%%%%%%%%%%%%%%%%%%%%%%%%%%
\begin{frame}[fragile]\frametitle{The 4Ds: Diligence}
\begin{columns}
    \begin{column}[T]{0.6\linewidth}
      \begin{itemize}
        \item \textbf{Diligence}: Using AI responsibly and ethically
        \item Three dimensions:
          \begin{itemize}
            \item \textbf{Creation Diligence:} Be thoughtful about \emph{which} AI systems you use and how you interact with them
            \item \textbf{Transparency Diligence:} Be honest about AI's role in your work with everyone who needs to know
            \item \textbf{Deployment Diligence:} Take responsibility for verifying and vouching for the outputs you use or share
          \end{itemize}
        \item Diligence means you remain \emph{accountable} for AI-assisted work: the AI is a tool, not the decision-maker
      \end{itemize}
    \end{column}
    \begin{column}[T]{0.4\linewidth}
      \begin{alertblock}{Key Question}
        \emph{``Am I using AI in a way I would be comfortable disclosing and defending?''}
      \end{alertblock}
      \vspace{0.5em}
      \begin{exampleblock}{Example}
        Disclosing AI assistance in a report, and verifying all facts before publication
      \end{exampleblock}
    \end{column}
  \end{columns}
\end{frame}

%%%%%%%%%%%%%%%%%%%%%%%%%%%%%%%%%%%%%%%%%%%%%%%%%%%%%%%%%%%%%%%%%%%%%%%%%%%%%%%%%%
\begin{frame}[fragile]\frametitle{Human--AI Interaction Modes}
\begin{columns}
    \begin{column}[T]{0.55\linewidth}
      \begin{itemize}
        \item \textbf{Automation}
          \begin{itemize}
            \item AI performs specific tasks based on specific human instructions
            \item Human defines \emph{what} needs to be done; AI executes it
            \item Example: generating a formatted report from structured data
          \end{itemize}
        \item \textbf{Augmentation}
          \begin{itemize}
            \item Humans and AI collaborate as thinking partners
            \item Iterative back-and-forth; both contribute to the outcome
            \item Example: co-writing and refining a document together
          \end{itemize}
        \item \textbf{Agency}
          \begin{itemize}
            \item AI works independently on the human's behalf
            \item Human sets goals and behaviour patterns, not exact actions
            \item Example: an AI agent that schedules meetings autonomously
          \end{itemize}
      \end{itemize}
    \end{column}
    \begin{column}[T]{0.45\linewidth}
      \begin{block}{Spectrum of Control}
        \begin{center}
          \begin{tabular}{cl}
            \textbf{Mode} & \textbf{Human role} \\
            \hline
            Automation  & Director \\
            Augmentation & Collaborator \\
            Agency      & Goal-setter \\
          \end{tabular}
        \end{center}
      \end{block}
      \vspace{0.4em}
      Choosing the right mode is itself a \textbf{Delegation} skill.
    \end{column}
  \end{columns}
\end{frame}

%%%%%%%%%%%%%%%%%%%%%%%%%%%%%%%%%%%%%%%%%%%%%%%%%%%%%%%%%%%%%%%%%%%%%%%%%%%%%%%%%%
\begin{frame}[fragile]\frametitle{AI Technical Concepts (1/2)}
% \begin{columns}
    % \begin{column}[T]{0.5\linewidth}
      \begin{itemize}
        \item \textbf{Generative AI}: AI systems that \emph{create} new content (text, images, code) rather than merely analysing existing data
        \item \textbf{Large Language Models (LLMs)}: Generative AI trained on vast text corpora to understand and generate human language; e.g.\ Claude
        \item \textbf{Parameters}: Mathematical values within a model that determine how it processes information; modern LLMs contain \emph{billions}
        \item \textbf{Neural Networks}: Interconnected nodes in layers that learn patterns from data; inspired by, but distinct from, biological brains
        \item \textbf{Transformer Architecture}: Breakthrough 2017 design enabling LLMs to process text in parallel while tracking relationships across long passages
      % \end{itemize}
    % \end{column}
    % \begin{column}[T]{0.5\linewidth}
      % \begin{itemize}
        \item \textbf{Pre-training}: Initial phase where a model learns patterns from massive text datasets, building foundational language understanding
        \item \textbf{Fine-tuning}: Additional training to make a model follow instructions, give helpful responses, and avoid harmful content
        \item \textbf{Scaling Laws}: Empirical observation that as models grow larger and train on more data, performance improves consistently; new capabilities can \emph{emerge} at scale thresholds
        \item \textbf{Reasoning / Thinking Models}: Models designed to work step-by-step through complex problems, showing improved logical reasoning
      \end{itemize}
    % \end{column}
  % \end{columns}
\end{frame}

%%%%%%%%%%%%%%%%%%%%%%%%%%%%%%%%%%%%%%%%%%%%%%%%%%%%%%%%%%%%%%%%%%%%%%%%%%%%%%%%%%
\begin{frame}[fragile]\frametitle{AI Technical Concepts (2/2)}
% \begin{columns}
    % \begin{column}[T]{0.5\linewidth}
      \begin{itemize}
        \item \textbf{Context Window}: The amount of information an AI can consider at one time (conversation history + documents); has a maximum limit that varies by model
        \item \textbf{Hallucination}: An error where the AI confidently states something that sounds plausible but is actually incorrect
        \item \textbf{Knowledge Cutoff Date}: The point after which the model has no built-in knowledge, based on when training data ended
        \item \textbf{Temperature}: A setting controlling output randomness: higher temperature $\rightarrow$ more creative; lower temperature $\rightarrow$ more predictable
      % \end{itemize}
    % \end{column}
    % \begin{column}[T]{0.5\linewidth}
      % \begin{itemize}
        \item \textbf{Retrieval Augmented Generation (RAG)}: Connects AI to external knowledge sources at inference time, improving accuracy and reducing hallucinations
        \item \textbf{Bias}: Systematic patterns in AI outputs that unfairly favour or disadvantage certain groups, often reflecting imbalances in training data
      \end{itemize}
      \vspace{0.4em}
      \begin{alertblock}{Practical implication}
        Hallucination + knowledge cutoff + bias together mean \textbf{Discernment} is always necessary: never accept AI output uncritically.
      \end{alertblock}
    % \end{column}
  % \end{columns}
\end{frame}

%%%%%%%%%%%%%%%%%%%%%%%%%%%%%%%%%%%%%%%%%%%%%%%%%%%%%%%%%%%%%%%%%%%%%%%%%%%%%%%%%%
\begin{frame}[fragile]\frametitle{Prompt Engineering Concepts}
% \begin{columns}
    % \begin{column}[T]{0.5\linewidth}
      \begin{itemize}
        \item \textbf{Prompt}: The full input given to an AI model: instructions, context, documents, and examples
        \item \textbf{Prompt Engineering}: The practice of designing effective prompts to produce desired outputs; combines clear communication with AI-specific techniques
        \item \textbf{Chain-of-Thought Prompting}: Encouraging the AI to work step-by-step through a problem before giving a final answer; improves accuracy on complex tasks
        \item \textbf{Few-Shot Learning (n-shot prompting)}: Providing $n$ input--output examples so the model understands the desired pattern without lengthy explanations
      % \end{itemize}
    % \end{column}
    % \begin{column}[T]{0.5\linewidth}
      % \begin{itemize}
        \item \textbf{Role / Persona Definition}: Specifying a character, expertise level, or communication style for the AI to adopt (e.g.\ ``Explain this as a UX expert'')
        \item \textbf{Output Constraints / Formatting}: Specifying the desired format, length, or structure within the prompt% (e.g.\ ``Reply in a 3-column table, max 200 words'')
        \item \textbf{Think-First Approach}: Explicitly asking the AI to reason through its thinking \emph{before} providing a final answer, leading to more considered responses
      \end{itemize}
    % \end{column}
  % \end{columns}
\end{frame}

%%%%%%%%%%%%%%%%%%%%%%%%%%%%%%%%%%%%%%%%%%%%%%%%%%%%%%%%%%%%%%%%%%%%%%%%%%%%%%%%%%
\begin{frame}[fragile]\frametitle{Putting It All Together}
\begin{columns}
    \begin{column}[T]{0.6\linewidth}
      \begin{itemize}
        \item AI Fluency is not a single skill: it is a \textbf{framework} that combines technical understanding, communication, critical thinking, and ethical responsibility
        \item The \textbf{4Ds} operate as a cycle:
          \begin{enumerate}
            \item \textit{Delegate} the right tasks to AI
            \item \textit{Describe} your needs clearly
            \item \textit{Discern} the quality of outputs
            \item \textit{Diligence}: take responsibility for what you share
          \end{enumerate}
        \item Understanding technical concepts (LLMs, hallucination, context windows) helps you apply the 4Ds more effectively
        \item Prompt engineering skills make your \textbf{Description} more powerful
      \end{itemize}
    \end{column}
    \begin{column}[T]{0.4\linewidth}
      \begin{block}{Remember}
        \begin{itemize}
          \item AI is a \emph{tool}, not an authority
          \item You remain \emph{accountable} for AI-assisted outputs
          \item Fluency grows with \emph{practice} and reflection
        \end{itemize}
      \end{block}
      \vspace{0.5em}
      {\tiny Source: Rick Dakan, Joseph Feller \& Anthropic, \textit{AI Fluency: Key Terminology Cheat Sheet}, 2025. CC BY-NC-SA 4.0}
    \end{column}
  \end{columns}
\end{frame}

%%%%%%%%%%%%%%%%%%%%%%%%%%%%%%%%%%%%%%%%%%%%%%%%%%%%%%%%%%%%%%%%%%%%%%%%%%%%%%%%%%